\documentclass[a4paper,11pt]{jsarticle}
\usepackage[dvipdfmx]{graphicx}
\usepackage{geometry}
\usepackage{longtable}
\usepackage{amsmath}
\geometry{margin=22mm}

\title{p2MEG run8データに対する時間相関解析の検証レポート}
\author{p2MEG 解析メモ（Codex支援）}
\date{2026年2月23日}

\begin{document}
\maketitle

\section{概要}
この文書は、\texttt{data/rawdata/mainexp/120$^\circ$/run8} の測定データについて、
次の2点を検証した結果をまとめたものである。

\begin{itemize}
  \item 時間差ヒストグラムに現れる3つの山が、どの条件で現れるか。
  \item エネルギーヒストグラムが「低エネルギー集中」または「山形」に見える条件は何か。
\end{itemize}

本レポートは、
\texttt{check/study\_michel\_ps\_tagged.C} と
\texttt{check/hist\_NaI\_fitint\_clean.C} の結果を比較し、
同じ再構成手順で条件だけを段階的に変える追加検証
（\texttt{check/report\_michel\_ps\_analysis.C}）に基づいて作成した。

\section{用語と記号の定義}
この節で定義した用語だけを、以降で使用する。

\begin{longtable}{p{0.25\linewidth}p{0.70\linewidth}}
プラスチックシンチレータ（PS） & 陽電子の通過を速い時刻情報として与える検出器。PS\_A, PS\_B の2チャネルを使う。 \\
ヨウ化ナトリウムシンチレータ（NaI） & エネルギー情報を与える検出器。NaI\_A1, NaI\_A2, NaI\_B1, NaI\_B2 の4チャネルを使う。 \\
デジタイザ事象 & 波形を1000サンプルずつに区切った1単位。 \\
サンプル & 波形の時間ビン。1サンプルは4 ns。 \\
モジュールA, モジュールB & モジュールAは NaI\_A1 と NaI\_A2 の組、モジュールBは NaI\_B1 と NaI\_B2 の組。 \\
$ t_{\mathrm{PS}} $ & PSの代表時刻。標準定義は「PS候補のうち最も早い時刻」。 \\
$ t_{\mathrm{mod}} $ & NaIモジュール時刻。対応した2チャネル時刻の平均。 \\
$ \Delta t $ & 時間差。$\Delta t = t_{\mathrm{mod}} - t_{\mathrm{PS}}$。単位はサンプル。 \\
即時領域 & $-20 \le \Delta t \le 20$ を指す解析上の呼称。 \\
遅延領域 & $110 \le \Delta t \le 170$ を指す解析上の呼称。 \\
負側領域 & $-220 \le \Delta t \le -100$ を指す解析上の呼称。 \\
$N(E_1,E_2)$ & エネルギー区間 $[E_1,E_2]$ に入る事象数。 \\
$f_{\mathrm{low}}$ & 低エネルギー比率。$f_{\mathrm{low}} = N(0,20000) / N(0,200000)$。 \\
\end{longtable}

\section{解析手順}
\subsection{入力と再構成}
入力は run8 の6波形ファイル
（PS: 2チャネル、NaI: 4チャネル）である。
6チャネルを同じ時刻順で同時に読み込み、1000サンプルを1デジタイザ事象として処理した。
時刻単位は1サンプル = 4 nsである。

本解析は、次の2段階で実行した。

\paragraph{Pass 1: PSしきい値のデータ駆動決定}
PSしきい値を固定値で与えず、run8データから自動決定した。
各デジタイザ事象で以下を行う。
\begin{enumerate}
  \item 波形のベースラインを、その事象内のADC最頻値で求める。
  \item パルス波形を $s_i = \mathrm{baseline} - \mathrm{ADC}_i$ で定義する。
  \item 単純パルス抽出条件（立上がりしきい値、終了しきい値、最小分離）でPS候補を抽出し、振幅ヒストグラムを作る。
  \item 振幅ヒストグラム上で、主ピークより低振幅側にある谷を自動探索し、その谷位置を採用しきい値とする。
\end{enumerate}

\paragraph{Pass 2: 事象ごとの時刻再構成とエネルギー再構成}
Pass 1で得たPSしきい値を固定して、各事象で本再構成を行う。
\begin{enumerate}
  \item PS候補のうち、振幅がしきい値以上の候補だけを残す。
  残った候補の中で最も早い時刻を $t_{\mathrm{PS}}$ と定義する。
  \item NaI 4チャネルそれぞれで、パルス分解付きフィットを行い、各パルスの時刻 $t_0$ と積分値 $I$ を求める。
  フィット形は
  \(
    s(t)=A\{\exp[-(t-t_0)/(t_r+d)]-\exp[-(t-t_0)/t_r]\}
  \)
  （$t<t_0$ では0）とし、積分値は $I=A\cdot d$ とした。
  \item NaI\_A1とNaI\_A2（同様にNaI\_B1とNaI\_B2）を時刻差30サンプル以内で対応付ける。
  対応付けた1組に対して、モジュール時刻を2チャネル時刻の平均、モジュール積分値を2チャネル積分値の和で定義する。
  \item 各モジュール集合の中から、$t_{\mathrm{PS}}$ に最も近い1組を採用する。
  このとき時間差
  \(
    \Delta t = t_{\mathrm{mod}} - t_{\mathrm{PS}}
  \)
  を計算し、時間差ヒストグラムへ加える。
\end{enumerate}

この手順により、
「PSで時刻を定義した同時事象」と
「NaIで再構成したエネルギー」の対応を一貫した規則で作った。

\subsection{しきい値と探索窓の設定根拠}
\paragraph{PS振幅しきい値}
固定値は使わず、PS振幅ヒストグラムの谷を自動探索して決めた。
得られた値は
\(\mathrm{thr}_{\mathrm{PS,A}}=125.5\)、
\(\mathrm{thr}_{\mathrm{PS,B}}=178.5\)（ADCカウント）である。
これは「低振幅成分」と「主ピーク」の境界をデータ駆動で決めるためである。

図\ref{fig:ps-threshold}に、PS振幅ヒストグラムと採用しきい値を示す。
黒破線は採用しきい値、緑破線は谷位置、紫破線は主ピーク位置である。
しきい値以上に残る割合は、
PS\_Aで36.3\%（32175/88685）、
PS\_Bで51.4\%（3047/5924）であった。
低振幅側を除去しつつ、主ピーク側を保持していることが確認できる。

\begin{figure}[htbp]
  \centering
  \includegraphics[width=0.95\linewidth]{mainexp/plot_ps_threshold_validation_run8.pdf}
  \caption{PS振幅ヒストグラムと自動決定しきい値。}
  \label{fig:ps-threshold}
\end{figure}

\paragraph{時間差ピーク探索窓}
時間差分布の主ピークを分離するため、
即時側の探索窓を $[-60,60]$、遅延側の探索窓を $[80,260]$ とした。
これは解析上の運用窓であり、物理定数ではない。

\section{結果1: 時間差分布の形}
図\ref{fig:dt}に、標準定義（最も早いPS時刻を基準）で作った
時間差 $\Delta t$ 分布を示す。

\begin{figure}[htbp]
  \centering
  \includegraphics[width=0.95\linewidth,page=1]{mainexp/report_michel_ps_analysis_run8.pdf}
  \caption{時間差 $\Delta t$ 分布（標準定義）。}
  \label{fig:dt}
\end{figure}

ピーク位置は次のとおりである。
\begin{itemize}
  \item モジュールA: 即時側 1、遅延側 129（サンプル）
  \item モジュールB: 即時側 3、遅延側 129（サンプル）
\end{itemize}

領域別事象数は次のとおりである。
\begin{itemize}
  \item モジュールA: 負側 1618、即時 3208、遅延 1580
  \item モジュールB: 負側 311、即時 414、遅延 288
\end{itemize}

ここで、選別後PSの「2本目時刻−1本目時刻」分布を作ると、
ピークは 139 サンプルであった（図\ref{fig:ps-gap}）。
この量は、同一事象の中でPS信号が繰り返し現れる代表的な時間間隔を表す。

一方、時間差分布の遅延側ピークは 129 サンプルである。
この2つが近いという事実は、
「最も早いPS時刻を基準にしたとき、NaI側の対応パルスが次の時間群に属すると、
時間差が“1つ分の時間間隔”だけ正方向へずれる」事象が
遅延側ピークを作っている可能性が高いことを意味する。

差は 10 サンプル（40 ns）あるが、これは
時刻抽出のばらつき、NaI時刻を2チャネル平均で定義していること、
およびピーク位置をヒストグラム最大ビンで読む離散化の影響で説明可能な範囲である。

\begin{figure}[htbp]
  \centering
  \includegraphics[width=0.8\linewidth,page=3]{mainexp/report_michel_ps_analysis_run8.pdf}
  \caption{選別後PSの「2本目時刻−1本目時刻」分布。ピークは139サンプル。}
  \label{fig:ps-gap}
\end{figure}

\section{結果2: 3つの山の原因切り分け}
\subsection{PS基準時刻の定義を変えた比較}
標準定義（最も早いPS時刻）を、
「同一事象内のPS候補のうち、モジュール時刻との差が最小になる時刻」へ変更すると、
遅延領域が減り、即時領域が増加した（図\ref{fig:dtref}）。

\begin{figure}[htbp]
  \centering
  \includegraphics[width=0.95\linewidth,page=2]{mainexp/report_michel_ps_analysis_run8.pdf}
  \caption{PS基準時刻の定義比較（青: 標準定義、赤: 最小時間差定義）。}
  \label{fig:dtref}
\end{figure}

モジュールAでの変化は次のとおりである。
\begin{itemize}
  \item 遅延領域: 1580 $\rightarrow$ 1235
  \item 即時領域: 3208 $\rightarrow$ 3645
\end{itemize}

\subsection{負側領域の詳細検証}
モジュールAの負側領域（1608事象）について、
同じモジュール時刻を別のPS時刻へ再対応した。

\begin{itemize}
  \item PSしきい値を満たした候補のみを使う場合:
  即時領域へ戻る事象は 0（0\%）
  \item しきい値未満も含むPS候補を使う場合:
  即時領域へ戻る事象は 1082（67.3\%）
\end{itemize}

同様に遅延領域（1561事象）では、
しきい値未満も含めると 1094 事象（70.1\%）が即時領域へ戻る。

この結果は、負側領域と遅延領域の一部が
「PSしきい値によって代表時刻がずれる効果」で説明できることを示す。
ただし、負側領域には約33\%の残差があり、
偶発一致や波形分解誤差など他要因の寄与は残る。

\section{結果3: エネルギー分布形状の条件依存}
図\ref{fig:ecomp}に、モジュールAのエネルギー分布比較を示す。

\begin{figure}[htbp]
  \centering
  \includegraphics[width=0.95\linewidth,page=4]{mainexp/report_michel_ps_analysis_run8.pdf}
  \caption{モジュールAのエネルギー分布。左: 絶対事象数、右: 正規化形状。}
  \label{fig:ecomp}
\end{figure}

低エネルギー比率 $f_{\mathrm{low}}$ は次のとおりである。
\begin{itemize}
  \item NaI\_A1単独（PS条件なし）: 0.475
  \item A1+A2対応ペア和（PS条件なし）: 0.111
  \item A1+A2対応ペア和（PS条件あり）: 0.109
  \item A1+A2でPS最小時間差1組のみ: 0.109
  \item 上記に即時領域条件を追加: 0.0787
\end{itemize}

したがって、分布形状の主変化は
「PS条件の有無」単独ではなく、
\begin{enumerate}
  \item 2チャネル対応ペア和を取ること
  \item さらに時間差条件（最小時間差対応、即時領域選択）を課すこと
\end{enumerate}
によって生じる。

\subsection{対応ペア和で低エネルギー比率が減る理由}
「なぜ A1+A2 の対応ペア和を取ると低エネルギー比率が下がるか」は、
次の3つの効果で説明できる。

\paragraph{(1) 2チャネル同時成立条件による低エネルギー偶発の抑制}
1チャネルだけで見ると、低エネルギー側には
小振幅ノイズ、基線揺らぎ、片側だけで立った弱いパルスが多く含まれる。
対応ペア和では「A1とA2の両方に、時刻が近いパルスがあること」を要求するため、
この種の片側事象が大きく減る。

\paragraph{(2) 和を取ることで有効しきい値が上がる}
対応後の量は $I_{A1}+I_{A2}$ である。
2チャネルの和を使うため、単一チャネル分布に比べて
分布全体が高エネルギー側へ移動しやすい。
結果として、低エネルギー区間の占有率が下がる。

\paragraph{(3) 時間相関条件が追加で偶発を落とす}
さらに「PS基準で最小時間差の1組のみ」や「即時領域」条件を加えると、
時間的に無関係な組合せが減るため、低エネルギー側がもう一段減る。

実測値もこの流れと整合する。
\begin{itemize}
  \item NaI\_A1単独: $f_{\\mathrm{low}}=0.475$
  \item A1+A2対応ペア和（PS条件なし）: $f_{\\mathrm{low}}=0.111$
  \item A1+A2対応ペア和 + 即時領域: $f_{\\mathrm{low}}=0.0787$
\end{itemize}

なお、A1単独とA2単独の低エネルギー比率はそれぞれ 0.475 と 0.517 である。
両チャネルが完全に独立なら低エネルギー同時成立率は概算で $0.475\\times0.517\\approx0.246$ となるが、
実測の対応ペア和では 0.111 とさらに小さい。
これは、実際の対応条件（時刻一致、フィット品質、重なり除去）が
独立近似より強く低エネルギー偶発を抑制していることを示している。

\section{物理的考察}
\begin{itemize}
  \item 遅延側ピークは、PS時刻間隔ピークと近いことから、
  ビーム時間構造と事象内複数パルスの影響を強く受けている可能性が高い。
  \item 負側ピークは、PSしきい値起因の代表時刻ずれが大きく寄与する。
  ただし、全てをそれだけで説明することはできない。
  \item エネルギー分布については、
  1チャネル積分では低エネルギー事象が多いが、
  2チャネル同時成立と時間差条件を課すことで低エネルギー側が相対的に減少する。
\end{itemize}

\section{限界と今後の課題}
\begin{itemize}
  \item 本解析は1つのrun（120$^\circ$/run8）に基づく。
  他角度runでも同傾向かを検証する必要がある。
  \item 負側領域の残差要因（偶発一致、NaI時刻先行、フィット誤同定）の分離は未完了である。
  \item PSしきい値を連続的に変化させた系統評価を追加すると、
  代表時刻ずれの定量評価がより明確になる。
\end{itemize}

\section*{出力ファイル}
図を埋め込んだ完成版レポートは次に保存した。
\begin{itemize}
  \item \texttt{doc/report\_michel\_ps\_tagged\_run8.pdf}
\end{itemize}

\end{document}
